
\begin{claim}{D.26}{Proved; number-field strengthening of D.25}\label{clm:D26}
Let $k$ be any number field. If
$A\in\widehat{k^\times}=\varprojlim_n k^\times/k^{\times n}$
is an $n$th power in $k_v$ at all but finitely many finite
places $v$, for every $n\ge1$, then $A=1$.
The exceptional finite set may depend on $n$.

Consequently, every noncuspidal continuous section of
$\mathbf P^1_k-\{0,1,\infty\}$, and of every geometrically
connected finite \'etale cover of this curve, has local points
integral at a set of finite places of Dirichlet density one.
They are asymptotically integral in the precise sense of D.25.
No birational lift or cofinite integrality is asserted.
\end{claim}
\begin{proof}
\emph{Cyclotomic reduction of a finite Kummer component.}
Fix a prime power $n=\ell^r$, represent the component of $A$
by $b\in k^\times$, and put
$F=k(\mu_n,b^{1/n})$.
If $\operatorname{Gal}(F/k(\mu_n))$ had a nonidentity element
$\tau$, it would act on the roots of $X^n-b$ by a nontrivial
translation and fix none. Its conjugates have the same
fixed-point property. Chebotarev gives infinitely many finite
places, unramified in $F$ and away from $n$ and the divisor of
$b$, with this Frobenius conjugacy class. At such a place
$X^n-b$ has no root in $k_v$, because a root in $k_v$ would
be fixed by its unramified local Frobenius. This contradicts
the hypothesis. Thus a chosen root $\alpha$ lies in $k(\mu_n)$.

Write $\Gamma=\operatorname{Gal}(k(\mu_n)/k)$. The Kummer
class is the inflation of
\[
 c(\sigma)=\frac{\sigma(\alpha)}{\alpha},\qquad
 c\in H^1(\Gamma,\mu_n).
 \tag{D.26.1}
\]
For each $\sigma\in\Gamma$, the same Chebotarev argument
shows that $\sigma$ fixes some root $\alpha\zeta$ of
$X^n-b$. Otherwise its Frobenius class would again give
infinitely many local failures. Hence
$c(\sigma)=\zeta/\sigma(\zeta)$, and the restriction of $c$
to the cyclic subgroup generated by $\sigma$ is a coboundary.
This conclusion holds for every cyclic subgroup of $\Gamma$.

\emph{The finite-level exceptional kernel has exponent at most two.}
If $\ell$ is odd, the group $\Gamma$ is a subgroup of the
cyclic group $(\mathbf Z/\ell^r\mathbf Z)^\times$.
Thus $\Gamma$ itself is cyclic, $c=0$, and
$b\in k^{\times\ell^r}$.
For $\ell=2$ and $r\ge3$, put
\[
 \Gamma_0=\Gamma\cap
       (1+4\mathbf Z/2^r\mathbf Z).
\]
The group $1+4\mathbf Z/2^r\mathbf Z$ is cyclic, and
$[\Gamma:\Gamma_0]\le2$. The restriction of $c$ to
$\Gamma_0$ vanishes. Inflation--restriction therefore puts
$c$ in the image of
\[
 H^1(\Gamma/\Gamma_0,\mu_{2^r}^{\Gamma_0})
       \longrightarrow H^1(\Gamma,\mu_{2^r}).
\]
Positive-degree cohomology of a finite group is killed by
its order, so $2c=0$. Equivalently,
$b^2\in k^{\times 2^r}$. For $r=1,2$ the group $\Gamma$
is cyclic and this exponent bound also holds.
Every odd-primary component of $A$ is consequently trivial,
and every $2$-primary component has square one. Chinese
remaindering in the power completion gives $A^2=1$.

\emph{Compatibility removes the exceptional torsion.}
The group $k^\times/\mu(k)$ is free abelian. Indeed, the
unit subgroup modulo roots of unity is finitely generated
free abelian by the unit theorem, the image of the divisor
map is a subgroup of the free abelian group on the finite
primes, and the resulting extension splits because its
quotient is free abelian. The power completion of a free
abelian group is torsion-free: its coordinate maps for a
basis embed it in a product of copies of
$\widehat{\mathbf Z}$. It follows that the torsion subgroup
of $\widehat{k^\times}$ is exactly $\mu(k)$. Its subgroup
killed by two is $\{1,-1\}$.

Choose $r$ so large that $\mu_{2^{r+1}}$ is not contained
in $k$. The roots of $X^{2^r}+1$ are exactly the primitive
$2^{r+1}$st roots of unity. Every nonidentity element of
$\operatorname{Gal}(k(\mu_{2^{r+1}})/k)$ fixes none of them:
it acts by a unit $u\ne1$ modulo $2^{r+1}$, and multiplying
$u-1$ by an odd exponent cannot give zero modulo $2^{r+1}$.
Chebotarev yields infinitely many finite places where $-1$
is not a $2^r$th power. Thus $A=-1$ is incompatible with
the hypothesis, and $A=1$.
This proves the first assertion without a single-level
local-to-global power assertion or a Grunwald--Wang input.

\emph{The density conclusion over $k$.}
We check explicitly the number-field substitutions in D.25.
Every finite completion $k_v$ is a finite extension of some
$\mathbf Q_p$, so D.6 and D.19 give the local point or cusp.
At real places the real section theorem gives the required
connected component, and complex places give no further
condition. Thus the section is Selmer in the sense used by
Stix. His Theorem 10 is stated over every number field and
applies to the same Legendre family
\cite[Theorem 10]{D_Stix2015}.

The full Legendre $2$-torsion is still rational over the
three-punctured line. The three geometric cusp monodromies
and their distinct reductions in
$\mathbf P^1(\mathbf F_2)$ are therefore exactly those of
D.25. In the triangular alternative, at a sufficiently
good finite place $v$ of odd residue characteristic,
\[
 \chi_2(\operatorname{Frob}_v)=N(v)>1,
\]
so $\chi_2$ and $\chi_2^{-1}$ are distinct on $G_{k_v}$.
The same unique cyclotomically normalized line determines
one of the three $k$-rational cusps at every such place,
and that cusp is independent of $v$. At an exact cusp,
the argument uses the nontrivial geometric inertia and its
normalization character, including when the tangent cocycle
makes the section's arithmetic inertia trivial. This is
the same locally split case checked in D.25.

Let $c$ be the resulting cusp and choose respectively
$f=1-t$, $t$, or $t/(t-1)$, so that $f$ is a unit on the
open curve, regular at $c$, with $f(c)=1$.
The class $s^*\kappa(f)$ lies in $\widehat{k^\times}$.
At an ordinary point degenerating into $c$ it is a principal
unit; at an exact cuspidal section its value is $1$,
independently of the tangent class. For each $n$, Hensel's
lemma away from the places dividing $n$ and the finite
model exception set makes its local class trivial at
almost every finite place. The first assertion just proved
gives $s^*\kappa(f)=1$.

Local Kummer injectivity over each $k_v$ then excludes every
finite-local interior point, since none of the three displayed
functions takes the value $1$ on the open curve. Stix's
Corollary 5, with the corrected finite-place Stoll argument
already checked in D.22, applies over every number field.
The finite support is the three-point boundary; no assertion
about a real point lying in that boundary is made. Hence the
triangular alternative is cuspidal. A noncuspidal section
must have density-one integral places, and asymptotic
integrality is the other conclusion of Theorem 10.

Finally let $X\to\mathbf P^1_k-\{0,1,\infty\}$ be a
geometrically connected finite \'etale cover. Intersecting a
cusp decomposition group with its open fundamental subgroup
shows, exactly as in D.25, that its noncuspidal sections
push forward to noncuspidal sections. Normalize
$\mathbf P^1_{\mathcal O_k[1/S]}$ in $k(X)$ and enlarge the
finite set $S$ to obtain a good marked model and a finite
map, \'etale above the model of the three-punctured line.
The finite valuative criterion lifts integral local points
integrally. Removing finitely many places does not change
Dirichlet density. At ordinary degenerate places the local
boundary multiplicities satisfy
$d_v(\text{base})=e_y d_v(X)$, with only finitely many
possible positive integers $e_y$. Divisibility downstairs
by $n\prod_y e_y$ implies divisibility upstairs by $n$.
This proves both assertions for the finite \'etale cover.
\end{proof}
